% GENERATED FILE. Edit canonical lessons, then run npm run generate:books.
% canonical-source-hash: fnv1a64:93f23e7392eb8b58
% canonical-lessons: HI-C509-hartal, HI-C509-istifa, HI-C509-canda, HI-C509-manmani, HI-C509-nara

\chapter{Trouble at the Factory}
\label{ch:hi-trouble-at-the-factory}
\hlchaptermodality{509}

\begin{chapteropening}
\textbf{By the end of this chapter:} \emph{I can talk about strikes, collections and leaving a job.}

Complete the last lesson of chapter 509: I can talk about strikes, collections and leaving a job.
\end{chapteropening}

\section[\texorpdfstring{\hi{हड़ताल} (haṛtāl)}{हड़ताल (haṛtāl)}]{\hi{हड़ताल} (haṛtāl) — a strike (stopping work)}

\label{lesson:HI-C509-hartal}



\begin{quote}
Before the new one: say the Hindi for admission (to a school or course), then the Hindi for a salary, pay (formal).
\end{quote}

\subsection*{You'll want to know: \hi{हड़ताल}}
\textbf{\hi{हड़ताल}} — \emph{haṛtāl} — \textquotedblleft{}a strike (stopping work)\textquotedblright{}.

\begin{quote}
\hi{कारख़ाने} \hi{के} \hi{मज़दूर} \hi{तीन} \hi{दिन} \hi{से} \hi{हड़ताल} \hi{पर} \hi{हैं।}
\end{quote}

\begin{quote}
\emph{Kārkhāne ke mazdūr tīn din se haṛtāl par haiṁ.}
\end{quote}

\begin{quote}
The factory workers have been on strike for three days.
\end{quote}

Feminine.

\begin{cousinweb}
\textbf{\hi{हड़ताल}} (\emph{haṛtāl}) is filed by CDIAL under Sanskrit \emph{haṭṭa}, \textquotedblleft{}a market, a shop\textquotedblright{}, with the Hindi sense \textquotedblleft{}closing shops in mourning or protest\textquotedblright{} (CDIAL 13944). Today it is any stopping of work.
\end{cousinweb}

\begin{grammarlens}[title={What you've built}]
One more word for describing what you lived and felt.
\end{grammarlens}

\subsection*{Your turn}
\begin{itemize}
\raggedright
  \item \emph{Say it:} \emph{haṛtāl}
  \item \emph{Say it:} \emph{haṛtāl}, once more
  \item \emph{Say it:} \emph{Kārkhāne ke mazdūr tīn din se haṛtāl par haiṁ.} — The factory workers have been on strike for three days.
  \item \emph{Recall:} say the Hindi for admission (to a school or course), then the Hindi for a salary, pay (formal), then say \emph{haṛtāl} again
\end{itemize}

\begin{tcolorbox}[breakable,colback=teal!4,colframe=teal!35!black,title={Before you move on}]
What does \hi{हड़ताल} mean? (\textquotedblleft{}A strike (stopping work)\textquotedblright{}.) Say it once more.
\end{tcolorbox}
\section[\texorpdfstring{\hi{इस्तीफ़ा} (istīfā)}{इस्तीफ़ा (istīfā)}]{\hi{इस्तीफ़ा} (istīfā) — a resignation}

\label{lesson:HI-C509-istifa}



\begin{quote}
Before the new one: say the Hindi for a salary, pay (formal), then the Hindi for a strike (stopping work).
\end{quote}

\subsection*{You'll want to know: \hi{इस्तीफ़ा}}
\textbf{\hi{इस्तीफ़ा}} — \emph{istīfā} — \textquotedblleft{}a resignation\textquotedblright{}.

\begin{quote}
\hi{मीरा} \hi{ने} \hi{अपना} \hi{इस्तीफ़ा} \hi{मालिक} \hi{को} \hi{दे} \hi{दिया।}
\end{quote}

\begin{quote}
\emph{Mīrā ne apnā istīfā mālik ko de diyā.}
\end{quote}

\begin{quote}
Mira handed her resignation to the owner.
\end{quote}

Masculine.

\begin{cousinweb}
\textbf{\hi{इस्तीफ़ा}} (\emph{istīfā}) is usually traced to Arabic \emph{istiʿfā}, \textquotedblleft{}asking to be let off\textquotedblright{}, through Persian.
\end{cousinweb}

\begin{grammarlens}[title={What you've built}]
One more word for describing what you lived and felt.
\end{grammarlens}

\subsection*{Your turn}
\begin{itemize}
\raggedright
  \item \emph{Say it:} \emph{istīfā}
  \item \emph{Say it:} \emph{istīfā}, once more
  \item \emph{Say it:} \emph{Mīrā ne apnā istīfā mālik ko de diyā.} — Mira handed her resignation to the owner.
  \item \emph{Recall:} say the Hindi for a salary, pay (formal), then the Hindi for a strike (stopping work), then say \emph{istīfā} again
\end{itemize}

\begin{tcolorbox}[breakable,colback=teal!4,colframe=teal!35!black,title={Before you move on}]
What does \hi{इस्तीफ़ा} mean? (\textquotedblleft{}A resignation\textquotedblright{}.) Say it once more.
\end{tcolorbox}
\section[\texorpdfstring{\hi{चंदा} (candā)}{चंदा (candā)}]{\hi{चंदा} (candā) — a contribution; a collection (of money)}

\label{lesson:HI-C509-canda}



\begin{quote}
Before the new one: say the Hindi for a strike (stopping work), then the Hindi for a resignation.
\end{quote}

\subsection*{You'll want to know: \hi{चंदा}}
\textbf{\hi{चंदा}} — \emph{candā} — \textquotedblleft{}a contribution; a collection (of money)\textquotedblright{}.

\begin{quote}
\hi{मज़दूरों} \hi{ने} \hi{एक} \hi{बीमार} \hi{साथी} \hi{के} \hi{लिए} \hi{चंदा} \hi{दिया।}
\end{quote}

\begin{quote}
\emph{Mazdūroṁ ne ek bīmār sāthī ke liye candā diyā.}
\end{quote}

\begin{quote}
The workers chipped in money for a sick companion.
\end{quote}

Masculine.

\begin{cousinweb}
\textbf{\hi{चंदा}} (\emph{candā}) is usually traced to Persian \emph{chanda}, \textquotedblleft{}a contribution\textquotedblright{}, built on \emph{chand}, \textquotedblleft{}how many, how much\textquotedblright{}: what each person puts in. It is not the \emph{candā}, \textquotedblleft{}moon\textquotedblright{}, that CDIAL derives from Sanskrit \emph{candra} (CDIAL 4661).
\end{cousinweb}

\begin{grammarlens}[title={What you've built}]
One more word for describing what you lived and felt.
\end{grammarlens}

\subsection*{Your turn}
\begin{itemize}
\raggedright
  \item \emph{Say it:} \emph{candā}
  \item \emph{Say it:} \emph{candā}, once more
  \item \emph{Say it:} \emph{Mazdūroṁ ne ek bīmār sāthī ke liye candā diyā.} — The workers chipped in money for a sick companion.
  \item \emph{Recall:} say the Hindi for a strike (stopping work), then the Hindi for a resignation, then say \emph{candā} again
\end{itemize}

\begin{tcolorbox}[breakable,colback=teal!4,colframe=teal!35!black,title={Before you move on}]
What does \hi{चंदा} mean? (\textquotedblleft{}A contribution; a collection (of money)\textquotedblright{}.) Say it once more.
\end{tcolorbox}
\section[\texorpdfstring{\hi{मनमानी} (manmānī)}{मनमानी (manmānī)}]{\hi{मनमानी} (manmānī) — high-handedness; doing as one pleases}

\label{lesson:HI-C509-manmani}



\begin{quote}
Before the new one: say the Hindi for a resignation, then the Hindi for a contribution.
\end{quote}

\subsection*{You'll want to know: \hi{मनमानी}}
\textbf{\hi{मनमानी}} — \emph{manmānī} — \textquotedblleft{}high-handedness; doing as one pleases\textquotedblright{}.

\begin{quote}
\hi{मालिक} \hi{की} \hi{मनमानी} \hi{से} \hi{सब} \hi{मज़दूर} \hi{परेशान} \hi{हैं।}
\end{quote}

\begin{quote}
\emph{Mālik kī manmānī se sab mazdūr pareśān haiṁ.}
\end{quote}

\begin{quote}
All the workers are fed up with the owner's high-handedness.
\end{quote}

Feminine.

\begin{cousinweb}
\textbf{\hi{मनमानी}} (\emph{manmānī}) joins \emph{man}, \textquotedblleft{}mind\textquotedblright{}, from Sanskrit \emph{manas} (CDIAL 9822), and a form of \emph{mānnā}, \textquotedblleft{}to respect, obey\textquotedblright{}, from Sanskrit \emph{manyate}, \textquotedblleft{}thinks fit\textquotedblright{} (CDIAL 9857): doing only what one's own mind will have.
\end{cousinweb}

\begin{grammarlens}[title={What you've built}]
One more word for describing what you lived and felt.
\end{grammarlens}

\subsection*{Your turn}
\begin{itemize}
\raggedright
  \item \emph{Say it:} \emph{manmānī}
  \item \emph{Say it:} \emph{manmānī}, once more
  \item \emph{Say it:} \emph{Mālik kī manmānī se sab mazdūr pareśān haiṁ.} — All the workers are fed up with the owner's high-handedness.
  \item \emph{Recall:} say the Hindi for a resignation, then the Hindi for a contribution, then say \emph{manmānī} again
\end{itemize}

\begin{tcolorbox}[breakable,colback=teal!4,colframe=teal!35!black,title={Before you move on}]
What does \hi{मनमानी} mean? (\textquotedblleft{}High-handedness; doing as one pleases\textquotedblright{}.) Say it once more.
\end{tcolorbox}
\section[\texorpdfstring{\hi{नारा} (nārā)}{नारा (nārā)}]{\hi{नारा} (nārā) — a slogan; a rallying cry}

\label{lesson:HI-C509-nara}



\begin{quote}
Before the new one: say the Hindi for a contribution, then the Hindi for high-handedness.
\end{quote}

\subsection*{You'll want to know: \hi{नारा}}
\textbf{\hi{नारा}} — \emph{nārā} — \textquotedblleft{}a slogan; a rallying cry\textquotedblright{}.

\begin{quote}
\hi{मज़दूरों} \hi{के} \hi{नारे} \hi{कारख़ाने} \hi{के} \hi{अंदर} \hi{तक} \hi{सुनाई} \hi{दे} \hi{रहे} \hi{थे।}
\end{quote}

\begin{quote}
\emph{Mazdūroṁ ke nāre kārkhāne ke andar tak sunāī de rahe the.}
\end{quote}

\begin{quote}
The workers' slogans could be heard right inside the factory.
\end{quote}

Masculine.

\begin{cousinweb}
\textbf{\hi{नारा}} (\emph{nārā}) is usually traced to Arabic \emph{naʿra}, \textquotedblleft{}a shout, a cry\textquotedblright{}, through Persian; Hindi uses it for the words a crowd shouts together.
\end{cousinweb}

\begin{grammarlens}[title={What you've built}]
One more word for describing what you lived and felt.
\end{grammarlens}

\subsection*{Your turn}
\begin{itemize}
\raggedright
  \item \emph{Say it:} \emph{nārā}
  \item \emph{Say it:} \emph{nārā}, once more
  \item \emph{Say it:} \emph{Mazdūroṁ ke nāre kārkhāne ke andar tak sunāī de rahe the.} — The workers' slogans could be heard right inside the factory.
  \item \emph{Recall:} say the Hindi for a contribution, then the Hindi for high-handedness, then say \emph{nārā} again
\end{itemize}

\begin{tcolorbox}[breakable,colback=teal!4,colframe=teal!35!black,title={Before you move on}]
What does \hi{नारा} mean? (\textquotedblleft{}A slogan; a rallying cry\textquotedblright{}.) Say it once more.
\end{tcolorbox}
